\section{Feature selection on eight held-out data sources}
\label{sec:minpi-compression}

\begin{table}[t]
\centering
\caption{Held-out predictions on four classification and four regression sources,
with five splits per source. Width is the percentage of input columns retained.
Time in milliseconds includes ranking, width selection, final refitting and
prediction. Widths and times are geometric means, using the median of three
timings per split. Scores are arithmetic means with equal weight per source.
sMAE divides test MAE by the development response standard deviation.}
\label{tab:minpi-compression}
\small
\setlength{\tabcolsep}{4pt}
\begin{tabular}{@{}lrrrrrrrr@{}}
\toprule
& \multicolumn{4}{c}{Classification} & \multicolumn{4}{c}{Regression} \\
\cmidrule(lr){2-5}\cmidrule(l){6-9}
Ranking & Width & AUC & Accuracy & Time & Width & $R^2$ & sMAE & Time \\
\midrule
Full ridge       & 100.0 & .929 & .726 &  3.3 & 100.0 & .596 & .443 &  2.0 \\
Univariate $F$   &  56.4 & .922 & .709 & 61.2 &  46.1 & .568 & .459 &  9.5 \\
Ridge coefficient&  68.1 & .927 & .704 & 60.4 &  50.4 & .607 & .445 &  9.4 \\
Ridge RFE        &  69.1 & .926 & .707 & 68.8 &  47.0 & .562 & .462 & 12.8 \\
Mean utility gain&  65.0 & .920 & .722 & 68.9 &  47.4 & .594 & .439 & 13.9 \\
Sampled MinPI    &  68.9 & .924 & .713 & 68.6 &  50.4 & .602 & .446 & 13.1 \\
\bottomrule
\end{tabular}
\end{table}

Relative to full ridge, the primary MinPI rule met the prespecified regression
target. It retained 50.4\% of the columns, with a change of \(+.0066\) in
\(R^2\) and \(+.0033\) in scaled MAE (Table~\ref{tab:minpi-compression}).
The classification rule missed its accuracy target. AUC fell by .0048 and
accuracy fell by 1.29 percentage points against a one-point allowance
(95\% paired source-bootstrap interval for the accuracy change, \([-2.21,-0.38]\)
points). Resampling the four regression sources gave a retained-width interval
of \([25.1,90.5]\%\), spanning the 85\% target. These intervals describe
variation in the mix of the four observed sources in each domain.

Voting changed the selected columns on six of the 40 splits relative to
ranking by the full-model coefficients alone. It increased selection time by
14.9\% in classification and 50.9\% in regression. Against unit-step ridge
RFE, the classification width ratio was .998 (source-bootstrap interval
\([.870,1.112]\)). In regression, MinPI used 7.2\% more columns than RFE and
gained .041 in \(R^2\). The mean-gain control changes the thresholding,
aggregation and stopping rules together.

\paragraph{Sensitivity to the width rule.}
After the pilot and before scoring the eight sources, we specified a second
comparison requiring the selected width to match or improve on both full-model
validation scores. All six methods used this rule. MinPI retained 77.1\% of
columns in classification and 81.2\% in regression. Relative to full ridge,
classification AUC fell by .0024 and accuracy by 0.69 points. Regression
\(R^2\) rose by .0216 and scaled MAE fell by .0081. Both domains met the
prespecified loss and width targets at their point estimates. Four sources
retained all columns on every split. The regression difference from coefficient
ranking came from two Energy splits.

\paragraph{Negative witnesses and skipped work.}
Independent SVD refits over every coalition on the seven sources with at most
nine columns gave 47 exact negative indices in 245 feature--split pairs. The
sampler found 14 and labeled the other 33 zero. Five of the 14 witnessed
features survived width selection. Across all eight sources the sampler
scheduled 6,780 marginal evaluations. A previous negative witness saved 66
evaluations, and the losing-base bound saved 1,278. On the 64-column Leaves
data, every sampled index was zero and all 3,840 evaluations were required.
